\section{Related Work}
\label{sec:related}

\paragraph*{Attribution and faithfulness}
Attribution methods~\cite{simonyan2014saliency,sundararajan2017ig,lundberg2017shap} answer ``how relevant is unit $u$ to target $T$ under method $M$ and baseline $B$''. This answer can be insensitive to the model's parameters~\cite{adebayo2018sanity}, sensitive to function-preserving transformations~\cite{kindermans2019unreliability}, and different from attention weights~\cite{jain2019attention,wiegreffe2019notnot}. Jacovi and Goldberg separate plausibility from faithfulness~\cite{jacovi2020faithfulness}; removal tests such as ERASER comprehensiveness and sufficiency~\cite{deyoung2020eraser} and ROAR~\cite{hooker2019roar} operationalise faithfulness and depend on the replacement~\cite{hase2021ood}. In BeyondNN an attribution is \textsc{attributed} evidence and can never decide a causal relation.

\paragraph*{Interventions and causal abstraction}
Causal tracing and activation patching~\cite{meng2022rome,zhang2024patching} estimate how an output depends on a component. Causal scrubbing~\cite{chan2022scrubbing} turns a mechanistic hypothesis into a correspondence with a high-level graph tested by resampling ablations and summarised as recovered loss; causal abstraction~\cite{geiger2025abstraction} unifies patching, mediation and scrubbing, and distributed alignment search~\cite{geiger2024das} aligns high-level variables with representations. Causal scrubbing is the closest methodological precedent to BeyondNN's claim--test structure. BeyondNN differs in scope: it records evidence of several epistemic types from several protocols, keeps alternative interventions as named assumptions with roles, and returns a categorical standing with findings rather than a scalar.

\paragraph*{Probes and concepts}
Decodability does not imply use~\cite{belinkov2022probing}; control tasks~\cite{hewitt2019control} and amnesic probing~\cite{elazar2021amnesic} address this, concept activation vectors~\cite{kim2018tcav} measure sensitivity to a direction, and dictionary features~\cite{huben2024sae,bloom2024saelens} are often given generated labels whose text need not reflect the computation~\cite{turpin2023unfaithful}. BeyondNN keeps decodability, use and labels apart and derives \textsc{validated\_concept} only from passed encoding and use tests.

\paragraph*{Evaluation frameworks, benchmarks and statistics}
Quantus~\cite{hedstrom2023quantus} and OpenXAI~\cite{agarwal2022openxai} score explanation methods; Tracr~\cite{lindner2023tracr}, InterpBench~\cite{gupta2024interpbench}, MIB~\cite{mueller2025mib}, RAVEL~\cite{huang2024ravel} and CausalGym~\cite{arora2024causalgym} provide models or tasks with reference mechanisms. These score methods or proposed mechanisms; BeyondNN audits a particular claim under declared criteria and can itself be evaluated against their references (Section~\ref{sec:res-ref}). Several of these benchmarks define success through interchange interventions, the same family an audit uses, so agreement with them can reflect shared semantics. M\'eloux et al.~\cite{meloux2026variance} show that circuit-discovery estimates are high-variance; Certified Circuits~\cite{anani2026certified} certify the stability of circuit-inclusion decisions under data perturbation; and Certified Interventional Fidelity~\cite{asiaee2026cif} gives anytime-valid confidence sequences for interventional estimands. BeyondNN has none of these inferential guarantees; it records which estimand and assumptions underlie a standing.

\paragraph*{Claims and evidence}
Doshi-Velez and Kim ask that evaluation match the claim~\cite{doshivelez2017rigorous}; Joshi et al. use Pearl's hierarchy to state which claims each kind of evidence can justify~\cite{joshi2026causality}; Lin and Liu propose that causal claims disclose their identification strategy and stress-test its assumptions~\cite{lin2026identification}; and Lan et al. call for expert guidelines and source-based auditing of the arguments that interpretability claims depend on~\cite{lan2026auditable}. We do not claim that matching claims to evidence, or auditing interpretability arguments, is new. Our contribution is narrower: an implemented, machine-checkable semantics that matches typed evidence to scoped claims by provenance, applies declared \textsc{primary} and \textsc{alternative} rules to saved records, and retains counterevidence. Table~\ref{tab:capability} positions it against existing tools.

\paragraph*{Tooling}
Captum~\cite{kokhlikyan2020captum}, TransformerLens~\cite{nanda2022transformerlens}, NNsight~\cite{fiottokaufman2025nnsight}, pyvene~\cite{wu2024pyvene} and SAELens~\cite{bloom2024saelens} produce measurements and perform interventions. BeyondNN complements them: it consumes Captum attributions through an adapter, and its compiled-program study is cross-checked against an independent TransformerLens implementation.
